\section{JAM 2026}
\begin{enumerate}
\item For each positive integer $n$, let $x_n=1-(-1)^n+\frac1n$ and $y_n=\left(1+\frac1{2n}\right)^{3n}$. Which ONE of the following statements about the sequences $(x_n)$ and $(y_n)$ is TRUE?
\begin{enumerate}[label=(\Alph*)]
\item $(x_n)$ and $(y_n)$ are convergent.
\item $(x_n)$ is not convergent and $(y_n)$ is not convergent.
\item $(x_n)$ is convergent and $(y_n)$ is not convergent.
\item $(x_n)$ is not convergent and $(y_n)$ is convergent.
\end{enumerate}
\item For each positive integer $n$, define $x_n=(-1)^n$. Which ONE of the following statements about the sequence $(x_n)$ is FALSE?
\begin{enumerate}[label=(\Alph*)]
\item There exists $\epsilon>0$ such that $|x_n-1|<\epsilon$ for all positive integers $n$.
\item There exists $\epsilon>0$ such that for all $M>0$ there exists a positive integer $N>M$ for which $|x_N-1|>\epsilon$.
\item For all $\epsilon>0$ and $M>0$ there exists a positive integer $N$ such that $N>M$ and $|x_N-1|<\epsilon$.
\item For all $\epsilon>0$ and $M>0$ there exists a positive integer $N$ such that $N>M$ and $|x_N-1|>\epsilon$.
\end{enumerate}
\item Let $f:(1,4)\to\mathbb R$ be a differentiable function such that
\[f'(x)=(f(x)-\pi x)^2+\pi \qquad\text{for all }x\in(1,4).\]
Which ONE of the following is a possible value of $f(3)-f(2)$?
\begin{enumerate}[label=(\Alph*)]
\item $\pi+\frac16$
\item $\pi-\frac16$
\item $\frac{\pi}{2}+\frac13$
\item $\frac{\pi}{2}+\frac12$
\end{enumerate}
\item The general solution of the differential equation
\[x\sin\left(\frac yx\right)\frac{dy}{dx}=y\sin\left(\frac yx\right)+\frac x2,\qquad x>0\]
is given by \rule{2cm}{0.4pt}.
\begin{enumerate}[label=(\Alph*)]
\item $\cos(\frac xy)+\log_e(x^2)=K$, where $K$ is an arbitrary constant
\item $\cos(\frac yx)+\log_e(x^2)=K$, where $K$ is an arbitrary constant
\item $\cos(\frac xy)+\log_e(\sqrt x)=K$, where $K$ is an arbitrary constant
\item $\cos(\frac yx)+\log_e(\sqrt x)=K$, where $K$ is an arbitrary constant
\end{enumerate}
\item Let $f:\mathbb R^3\to\mathbb R$ be a continuous function. Then
\[\int_0^2\int_0^x\int_0^y f(x,y,z)\,dz\,dy\,dx=\rule{2cm}{0.4pt}.\]
\begin{enumerate}[label=(\Alph*)]
\item $\displaystyle\int_0^2\int_0^y\int_y^2 f(x,y,z)\,dx\,dz\,dy$
\item $\displaystyle\int_0^2\int_0^x\int_z^2 f(x,y,z)\,dy\,dz\,dx$
\item $\displaystyle\int_0^2\int_y^2\int_0^x f(x,y,z)\,dz\,dx\,dy$
\item $\displaystyle\int_0^2\int_z^2\int_0^x f(x,y,z)\,dy\,dx\,dz$
\end{enumerate}
\item Let $f:(1,2)\to\mathbb R$ be a continuous function satisfying
\[\lim_{x\to1^+}f(x)=3\quad\text{and}\quad\lim_{x\to2^-}f(x)=3.\]
Then which ONE of the following is necessarily TRUE?
\begin{enumerate}[label=(\Alph*)]
\item $f$ is bounded and $f$ has a maximum or a minimum but not both.
\item $f$ is bounded and $f$ has a maximum or a minimum or both.
\item $f$ is bounded and $f$ has neither a maximum nor a minimum.
\item $f$ is unbounded and $f$ has either no maximum or no minimum.
\end{enumerate}
\item Let $f:\mathbb R\to\mathbb R$ be a twice differentiable function such that $f(n)=1$ for all $n\in\mathbb Z$. Which ONE of the following statements is necessarily TRUE?
\begin{enumerate}[label=(\Alph*)]
\item $f'(n)=0$ for all $n\in\mathbb Z$.
\item $f'(x)=0$ for infinitely many $x\in\mathbb R\setminus\mathbb Z$.
\item $f''(n)=0$ for all $n\in\missingglyph$.
\item $f''(x)=0$ for infinitely many $x\in\mathbb R\setminus\missingglyph$.
\end{enumerate}
\item Let $A$ be an invertible $5\times5$ real matrix. Let $B$ be the matrix obtained by interchanging the second and third columns of $A$ and then adding $3$ times the third column to the fifth column. Which ONE of the following statements is TRUE?
\begin{enumerate}[label=(\Alph*)]
\item $B^{-1}$ is obtained by interchanging the second and third rows of $A^{-1}$, and then adding $3$ times the third row to the fifth row.
\item $B^{-1}$ is obtained by interchanging the second and third rows of $A^{-1}$, and then adding $-3$ times the fifth row to the third row.
\item $B^{-1}$ is obtained by adding $-3$ times the third row to the fifth row, and then interchanging the second and third rows of $A^{-1}$.
\item $B^{-1}$ is obtained by adding $3$ times the fifth row to the third row, and then interchanging the second and third rows of $A^{-1}$.
\end{enumerate}
\item Which ONE of the following statements is FALSE?
\begin{enumerate}[label=(\Alph*)]
\item $S_3$ is isomorphic to a subgroup of $S_4$.
\item $\missingglyph_3$ is isomorphic to a subgroup of $S_4$.
\item $S_3$ is isomorphic to a quotient group of $S_4$.
\item $\mathbb Z_6$ is isomorphic to a quotient group of $S_4$.
\end{enumerate}
\item Let $T$ be the triangular region in the plane with vertices at $(0,0),(1,0)$ and $(1,1)$. Let $f$ be a function defined from $T$ to $\mathbb R$ given in polar coordinates. Then the double integral of $f$ over $T$ is equal to \rule{2cm}{0.4pt}.
\begin{enumerate}[label=(\Alph*)]
\item $\displaystyle\int_0^{\sqrt2}\left(\int_0^{\pi/4}f\,d\theta\right)r\,dr+\int_1^{\sqrt2}\left(\int_0^{\cos^{-1}(1/r)}f\,d\theta\right)r\,dr$
\item $\displaystyle\int_0^{\sqrt2}\left(\int_0^{\pi/4}f\,d\theta\right)r\,dr-\int_1^{\sqrt2}\left(\int_0^{\cos^{-1}(1/r)}f\,d\theta\right)r\,dr$
\item $\displaystyle\int_0^{\sqrt2}\left(\int_0^{\pi/4}f\,d\theta\right)r\,dr+\int_1^{\sqrt2}\left(\int_0^{\sin^{-1}(1/r)}f\,d\theta\right)r\,dr$
\item $\displaystyle\int_0^{\sqrt2}\left(\int_0^{\pi/4}f\,d\theta\right)r\,dr-\int_1^{\sqrt2}\left(\int_0^{\sin^{-1}(1/r)}f\,d\theta\right)r\,dr$
\end{enumerate}
\item Given a sequence of real numbers $(a_n)$, define
\[b_n=\begin{cases}a_n,&a_n\ge0\\0,&a_n<0\end{cases}\quad\text{and}\quad
c_n=\begin{cases}a_n,&a_n<0\\0,&a_n\ge0\end{cases}\]
for each positive integer $n$. Which ONE of the following statements is TRUE?
\begin{enumerate}[label=(\Alph*)]
\item If $\sum_{n=1}^\infty a_n$ does not converge, then $\sum_{n=1}^\infty b_n$ does not converge.
\item If $\sum_{n=1}^\infty a_n$ converges but $\sum_{n=1}^\infty |a_n|$ does not converge, then $\sum_{n=1}^\infty b_n$ converges.
\item If $\sum_{n=1}^\infty a_n$ converges but $\sum_{n=1}^\infty |a_n|$ does not converge, then $\sum_{n=1}^\infty b_n$ does not converge.
\item If $\sum_{n=1}^\infty a_n$ converges, then $\sum_{n=1}^\infty b_n$ converges or $\sum_{n=1}^\infty c_n$ converges.
\end{enumerate}
\item For which one of the following pairs of values of $a$ and $b$, does the series
$\displaystyle\sum_{k=1}^\infty \frac{(3x)^k}{2^{\sqrt{k}}(1-x)^k}$
converge for all $x\in(a,b)$?
\begin{enumerate}[label=(\Alph*)]
\item $a=-\frac14$ and $b=\frac12$
\item $a=-1$ and $b=\frac15$
\item $a=-\frac12$ and $b=\frac14$
\item $a=-\frac14$ and $b=\frac13$
\end{enumerate}
\item Let $f:\mathbb R\to\mathbb R$ be a function that has derivatives of all orders. Which ONE of the following statements is FALSE?
\begin{enumerate}[label=(\Alph*)]
\item If $f$ is a polynomial in $x$ of degree at most $5$, then $f^{(6)}(x)=0$ for all $x\in\mathbb R$.
\item If $f^{(6)}(x)=0$ for all $x\in\mathbb R$, then $f$ is a polynomial in $x$ of degree at most $5$.
\item If $f^{(k)}(a)=0$ for all $1\le k\le5$ and $f^{(6)}(a)>0$ for some $a\in\mathbb R$, then $f(x)$ has a local minimum at $a$.
\item If $f^{(k)}(a)=0$ for all $1\le k\le6$ and $f^{(7)}(a)<0$ for some $a\in\mathbb R$, then $f(x)$ has a local maximum at $a$.
\end{enumerate}
\item Let $A$ be a $5\times3$ real matrix of rank $2$ and $b$ be a non-zero $5\times1$ real column vector. Suppose $\begin{pmatrix}1\\2\\3\end{pmatrix}$ and $\begin{pmatrix}4\\5\\6\end{pmatrix}$ are two solutions of the system of linear equations $Ax=b$. Which ONE of the following is also a solution of $Ax=b$?
\begin{enumerate}[label=(\Alph*)]
\item $\begin{pmatrix}3\\3\\3\end{pmatrix}$
\item $\begin{pmatrix}5\\7\\9\end{pmatrix}$
\item $\begin{pmatrix}5\\6\\9\end{pmatrix}$
\item $\begin{pmatrix}7\\8\\9\end{pmatrix}$
\end{enumerate}
\item Let $f:\mathbb R^2\to\mathbb R$ be the function given by
\[f(x,y)=(x^2-1)^2+(y^2-1)^2\quad\text{for all }(x,y)\in\mathbb R^2.\]
Which ONE of the following statements is TRUE?
\begin{enumerate}[label=(\Alph*)]
\item $f$ has local maxima at exactly two points.
\item $f$ has local minima at exactly two points.
\item $f$ has local minima at exactly three points.
\item $f$ has exactly four saddle points.
\end{enumerate}
\item The orthogonal trajectories of the family of curves
\[y=-3x-3+me^x\qquad\text{for a real parameter }m\]
are given by \rule{2cm}{0.4pt}.
\begin{enumerate}[label=(\Alph*)]
\item $x=\frac y3-\frac19+ke^{-3y}$, $k$ is a real parameter
\item $x=\frac y3-\frac19+ke^{3y}$, $k$ is a real parameter
\item $x=-\frac y3+\frac19+ke^{-3y}$, $k$ is a real parameter
\item $x=-\frac y3+\frac19+ke^{3y}$, $k$ is a real parameter
\end{enumerate}
\item Let $f:\mathbb R^2\to\mathbb R$ be the function defined by
\[f(x,y)=\begin{cases}\dfrac{5x^2y-4y^3}{x^2+y^2},&(x,y)\ne(0,0)\\0,&(\missingglyph,y)=(0,0).\end{cases}\]
Which ONE of the following statements is FALSE?
\begin{enumerate}[label=(\Alph*)]
\item $f$ is continuous at $(0,0)$
\item $\dfrac{\partial f}{\partial x}$ at $(0,1)=0$
\item $\dfrac{\partial f}{\partial y}$ at $(1,0)=5$
\item $\dfrac{\partial f}{\partial y}$ is continuous at $(0,0)$
\end{enumerate}
\item Let $\phi:(-\frac\pi2,\frac\pi2)\to\mathbb R$ be the solution of the differential equation
\[(\cos x)\frac{dy}{dx}-y=2y^2(\sin x-1)\cos x\]
satisfying $\phi(0)=\frac12$. Then $\phi(\frac\pi4)=\rule{2cm}{0.4pt}$.
\begin{enumerate}[label=(\Alph*)]
\item $2$
\item $\dfrac{3+2\sqrt2}{2}$
\item $\frac12$
\item $\dfrac1{\sqrt2}$
\end{enumerate}
\item The value of the following expression
\[\lim_{x\to\infty}\int_0^x e^{-t}|\sin t|\,dt\]
is \rule{2cm}{0.4pt}.
\begin{enumerate}[label=(\Alph*)]
\item $\dfrac{e^\pi+1}{2(e^\pi-1)}$
\item $\dfrac{e^\pi-1}{2(e^\pi+1)}$
\item $\dfrac{2(e^\pi+1)}{e^\pi-1}$
\item $\dfrac{2(e^\pi-1)}{e^\pi+1}$
\end{enumerate}
\item The semi-circle $y=\sqrt{6x-5-x^2}$ is rotated clockwise by the angle $\frac\pi2$ in the plane about the origin. What is the area of the planar region traced out by the semi-circle?
\begin{enumerate}[label=(\Alph*)]
\item $6\pi$
\item $4\pi$
\item $2\pi$
\item $\pi$
\end{enumerate}
\item Let $f:[1,9]\to\mathbb R$ be a non-constant continuous function such that $f(1)=f(9)$. Which ONE of the following statements is necessarily TRUE?
\begin{enumerate}[label=(\Alph*)]
\item There exists $c\in[4,8]$ such that $f(c)=f(c+1)$.
\item There exists $c\in[3,7]$ such that $f(c)=f(c+2)$.
\item There exists $c\in[2,6]$ such that $f(c)=f(c+3)$.
\item There exists $c\in[1,5]$ such that $f(c)=f(c+4)$.
\end{enumerate}
\item Let $S$ be the solid of intersection of two solid spheres $S_1,S_2$ given by
\[S_1:x^2+y^2+(z-1)^2\le4 \quad\text{and}\quad S_2:x^2+y^2+(z+1)^2\le4.\]
Which ONE of the following iterated integrals expresses the volume of $S$?
\begin{enumerate}[label=(\Alph*)]
\item $\displaystyle 8\int_0^{\sqrt3}\int_0^{\sqrt{3-x^2}}(\sqrt{4-x^2-y^2}-1)\,dy\,dx$
\item $\displaystyle 8\int_0^2\int_0^{\sqrt{4-x^2}}(\sqrt{4-x^2-y^2}-1)\,dy\,dx$
\item $\displaystyle 2\int_0^2\int_0^{\sqrt{4-x^2}}(\sqrt{4-x^2-y^2}-1)\,dy\,dx$
\item $\displaystyle 2\int_0^{\sqrt3}\int_0^{\sqrt{3-x^2}}(\sqrt{4-x^2-y^2}-1)\,dy\,dx$
\end{enumerate}
\item Consider the differential equation
\[\frac{d^2y}{dx^2}+3\frac{dy}{dx}+2y=\sin(e^x).\]
Which ONE of the following is a particular solution of this differential equation?
\begin{enumerate}[label=(\Alph*)]
\item $-e^{-x}\sin(e^x)$
\item $-e^{-2x}\sin(e^x)$
\item $-e^{-x}(\sin(e^x)+\cos(e^x))$
\item $-e^{-2x}(\sin(e^x)+\cos(e^x))$
\end{enumerate}
\item Let $e_1,e_2,e_3,e_4,e_5$ denote the standard basis vectors of the real vector space $\mathbb R^5$. Let $V$ be the subspace of $\mathbb R^5$ spanned by the vectors $e_1+e_2,e_2+e_3,e_3+e_4+e_5$. Consider the real vector space $M_{2,5}(\mathbb R)$. Let
\[S=\{P\in M_{2,5}(\mathbb R)\mid\operatorname{nullspace}(P)=V\}.\]
Which ONE of the following statements is TRUE?
\begin{enumerate}[label=(\Alph*)]
\item $S$ is the empty set.
\item $S=\{0\}$.
\item $S$ is non-empty and $S$ is not a subspace of $M_{2,5}(\mathbb R)$.
\item $S$ is a nonzero subspace of $M_{2,5}(\mathbb R)$.
\end{enumerate}
\item Let $P$ be a nonzero complex $n\times n$ matrix such that $v^*Pv\ge0$ for all column vectors $v$ in $\mathbb C^n$. Which ONE of the following statements is necessarily TRUE?
\begin{enumerate}[label=(\Alph*)]
\item All eigenvalues of $P$ are negative real numbers.
\item If $v_1,v_2$ are column vectors in $\mathbb C^n$ such that $Pv_1=v_2$ and $Pv_2=v_1$, then $v_1=v_2$.
\item $v^*Pv\ne0$ for all nonzero column vectors $v$ in $\mathbb C^n$.
\item All eigenvalues of $P^*$ are negative real numbers.
\end{enumerate}
\item Let $f:G\to\mathbb Z$ be a surjective group homomorphism from an abelian group $G$ to the additive group of integers $\mathbb Z$. Let $K$ denote the kernel of $f$. Which ONE of the following statements is FALSE?
\begin{enumerate}[label=(\Alph*)]
\item There exists a group homomorphism $\phi:\mathbb Z\to G$ such that $f\circ\phi$ is the identity homomorphism.
\item $K$ is isomorphic to a quotient group of $G$.
\item For all non-zero homomorphisms $\alpha:\mathbb Z\to G$ and $\beta:G\to K$, the composition $\beta\circ\alpha$ is non-zero.
\item $G$ has a subgroup isomorphic to $\mathbb Z$.
\end{enumerate}
\item Which ONE of the following statements is FALSE?
\begin{enumerate}[label=(\Alph*)]
\item There is a surjective group homomorphism from the additive group of rational numbers to the multiplicative group of all complex roots of unity.
\item The multiplicative group of complex numbers of modulus one is isomorphic to a quotient group of the additive group of real numbers.
\item Any group homomorphism from the multiplicative group of nonzero complex numbers into the group of all invertible $2\times2$ matrices with real entries has nontrivial kernel.
\item There exists a group homomorphism from the symmetric group on $n$ symbols into the multiplicative group of all invertible $n\times n$ matrices with real entries, which has a trivial kernel.
\end{enumerate}
\item For any two points $P,Q\in\mathbb R^3$, let $\operatorname{vect}(P,Q)$ denote the vector from the point $P$ to the point $Q$ and let $d(P,Q)$ denote the length of $\operatorname{vect}(P,Q)$. Let $F:\mathbb R^3\to\mathbb R^3$ be a linear transformation such that
\[d(F(P),F(Q))=d(P,Q)\quad\text{for all }P,Q\in\mathbb R^3.\]
Which ONE of the following statements is FALSE?
\begin{enumerate}[label=(\Alph*)]
\item $F$ is an injective map.
\item $F$ is a surjective map.
\item For any four points $P,Q,R,S$ in $\mathbb R^3$, we have
\[\operatorname{vect}(F(P),F(Q))\cdot\operatorname{vect}(F(R),F(S))=\operatorname{vect}(P,Q)\cdot\operatorname{vect}(R,S),\]
where $\cdot$ denotes the usual dot product in $\mathbb R^3$.
\item For any four points $P,Q,R,S$ in $\mathbb R^3$, we have
\[\operatorname{vect}(F(P),F(Q))\times\operatorname{vect}(F(R),F(S))=\operatorname{vect}(P,Q)\times\operatorname{vect}(R,S),\]
where $\times$ denotes the usual cross product in $\mathbb R^3$.
\end{enumerate}
\item For each positive integer $n$, let
\[x_n=1+\frac12+\cdots+\frac1n-\log_e n\]
and
\[y_n=\int_1^n\frac{\cos t}{t^2}\,dt.\]
Which ONE of the following statements about the sequences $(x_n)$ and $(y_n)$ is TRUE?
\begin{enumerate}[label=(\Alph*)]
\item $(x_n)$ and $(y_n)$ are convergent.
\item $(x_n)$ is convergent and $(y_n)$ is not convergent.
\item $(x_n)$ is not convergent and $(y_n)$ is convergent.
\item $(x_n)$ is not convergent and $(y_n)$ is not convergent.
\end{enumerate}
\item Let $G$ be the power set of $\{1,2,3,4\}$. Let the binary operation $\Delta$ on $G$ be
\[A\mathbin{\Delta} B=(A\setminus B)\cup(B\setminus A).\]
Which ONE of the following statements is TRUE?
\begin{enumerate}[label=(\Alph*)]
\item $\{1\}$ is an identity element of $(G,\Delta)$.
\item $(G,\Delta)$ is an abelian group but not cyclic.
\item $(G,\Delta)$ is a group and has an element of order $4$.
\item $(G,\Delta)$ is a group and has an element of order $8$.
\end{enumerate}
\item Which of the following functions $f:\mathbb R\to\mathbb R$ has/have a local minimum at $x=0$?
\begin{enumerate}[label=(\Alph*)]
\item $f(x)=\sin|x|$
\item $f(x)=\sin x+\dfrac{x^3}{6}$
\item $f(x)=x^4+x^2+3$
\item $f(x)=\min\{x-[x],1-x+[x]\}$
\end{enumerate}
\item Which of the following statements is/are TRUE?
\begin{enumerate}[label=(\Alph*)]
\item There exists a monotone sequence that does not converge but has a convergent subsequence.
\item There exists a sequence that has a bounded subsequence but does not have any convergent subsequence.
\item There exists a sequence $(x_n)$ such that given any positive integer $m$, $(x_n)$ has a subsequence converging to $m$.
\item There exists a sequence $(x_n)$ such that $(|x_{n+1}-x_n|)$ converges to $0$ but $(x_n)$ does not converge.
\end{enumerate}
\item Which of the following functions is/are differentiable at $x=1$?
\begin{enumerate}[label=(\Alph*)]
\item $f(x)=|x-1|^3$
\item $f(x)=|x^2-1|$
\item $f(x)=\begin{cases}x^2e^{-x^2},&|x|\le1\\e^{-1},&|x|>1.\end{cases}$
\item $f(x)=[x]$
\end{enumerate}
\item Let $f:\mathbb R\to\mathbb R$ be the function given by
\[f(x)=\begin{cases}\cos x,&x\in\mathbb Q\\0,&x\notin\mathbb Q.\end{cases}\]
Which of the following statements is/are TRUE?
\begin{enumerate}[label=(\Alph*)]
\item $f(x)$ is continuous at $0$.
\item $f(x)$ is continuous at $\frac\pi2$.
\item $f(x)$ is Riemann integrable on $[0,1]$ and $\int_0^1 f(x)\,dx=\sin1$.
\item $f(x)$ is Riemann integrable on $[0,1]$ and $\int_0^1 f(x)\,dx=0$.
\end{enumerate}
\item Consider the differential equation
\[(2y\cos x-xy\sin x)\,dx+2x\cos x\,dy=0\quad\text{for }x\in(0,\tfrac\pi4).\]
Which of the following is/are integrating factor(s) of the differential equation?
\begin{enumerate}[label=(\Alph*)]
\item $\dfrac1{xy}$
\item $xy$
\item $\sec x$
\item $\sqrt{\sec x}$
\end{enumerate}
\item Let $f(x,y)=\begin{cases}\dfrac{x^2y}{1+x^2}\sin(1/x),&x\ne0\\0,&x=0.\end{cases}$ Which of the following statements is/are TRUE?
\begin{enumerate}[label=(\Alph*)]
\item $\displaystyle\lim_{y\to0}\lim_{x\to0}f(x,y)$ exists.
\item $\dfrac{\partial f}{\partial x}$ is continuous at the point $(0,1)$.
\item $\dfrac{\partial f}{\partial y}$ is continuous at the point $(0,1)$.
\item $\displaystyle\lim_{(x,y)\to(0,0)}\frac{f(x,y)}{\sqrt{x^2+y^2}}$ does not exist.
\end{enumerate}
\item Let $P=\begin{bmatrix}0&1&0&0&0\\1&0&0&0&0\\0&0&0&1&0\\0&0&0&0&1\\0&0&1&0&0\end{bmatrix}\in M_5(\mathbb C)$. Which of the following statements is/are TRUE?
\begin{enumerate}[label=(\Alph*)]
\item $\operatorname{nullity}(P-I)\ge2$, where $I$ is the $5\times5$ identity matrix.
\item $P$ has $4$ distinct eigenvalues in $\mathbb C$.
\item $P$ has $3$ distinct eigenvalues in $\mathbb R$.
\item If $\lambda$ is an eigenvalue of $P$, then there exists a positive integer $n$ such that $\lambda^n=1$.
\end{enumerate}
\item Let $\binom nr$ denote the number of ways of choosing $r$ distinct objects out of $n$ distinct objects. Which of the following statements is/are TRUE?
\begin{enumerate}[label=(\Alph*)]
\item $\displaystyle\sum_{k=0}^6(-1)^k\binom{13}{2k}=-64$.
\item $\displaystyle\sum_{k=0}^6(-1)^k\binom{13}{2k+1}=128$.
\item $\displaystyle\sum_{k=0}^6(-1)^k\binom{13}{2k}=64$.
\item $\displaystyle\sum_{k=0}^6(-1)^k\binom{13}{2k+1}=-128$.
\end{enumerate}
\item Consider matrices $P$ of order $4\times6$ and $Q$ of the order $6\times4$ with real entries such that $PQ=0$ and $QP=0$. Which of the following statements is/are TRUE?
\begin{enumerate}[label=(\Alph*)]
\item $\operatorname{rangespace}(P)\subseteq\operatorname{nullspace}(Q)$ and $\operatorname{rangespace}(Q)\subseteq\operatorname{nullspace}(P)$.
\item $\operatorname{rank}(P)+\operatorname{rank}(Q)\le4$.
\item If $\operatorname{rangespace}(P)=\operatorname{nullspace}(Q)$, then $\operatorname{rank}(P)+\operatorname{rank}(Q)=4$.
\item $\operatorname{rangespace}(Q)=\operatorname{nullspace}(P)$.
\end{enumerate}
\item Let $V$ be the subset of $\mathbb R$ defined by
\[V=\left\{\frac{a+b\sqrt2}{c+d\sqrt2}:a,b,c,d\in\mathbb Q,\ c^2+d^2\ne0\right\}.\]
Which of the following statements is/are FALSE?
\begin{enumerate}[label=(\Alph*)]
\item $V$ is closed under the usual addition in $\mathbb R$.
\item $V$ is a subspace of the real vector space $\mathbb R$.
\item $V$ is a two dimensional subspace of the vector space $\mathbb R$ over $\mathbb Q$.
\item $V$ is a four dimensional subspace of the vector space $\mathbb R$ over $\mathbb Q$.
\end{enumerate}
\item The radius of convergence of the series $\displaystyle\sum_{n=1}^\infty\frac{(n!)^2}{(2n)!}(\log_e n)^{-1}x^n$ is \rule{2cm}{0.4pt} (rounded off to one decimal place).
\item $\displaystyle\int_0^1\left(\sum_{k=1}^\infty\frac{(\log_e2)^kx^{k-1}}{(k-1)!}\right)\,dx=$ \rule{2cm}{0.4pt} (rounded off to one decimal place).
\item $\displaystyle\lim_{n\to\infty}\left(\frac{n^2+1}{\sqrt{n^6+1}}+\cdots+\frac{n^2+n}{\sqrt{n^6+n}}\right)=$ \rule{2cm}{0.4pt} (rounded off to one decimal place).
\item $\displaystyle\lim_{x\to0}\left(\frac1x-\frac1{\sin x}+e^{(1-\cos x)/x}\right)=$ \rule{2cm}{0.4pt} (rounded off to one decimal place).
\item Let $\alpha$ and $\beta$ be real numbers such that the differential equation
\[(y^3+\alpha xy^4-5x+\cos 2y)\,dx+(3xy^2+20x^2y^3+\beta x\sin 2y)\,dy=0\]
is exact. Then $\alpha+\beta=$ \rule{2cm}{0.4pt} (rounded off to one decimal place).
\item Let $z=\cos(4x+5y)$, where $x=\frac\pi2+2\theta$, $y=-\left(\frac\pi4+\theta\right)$. Then $\left.\frac{dz}{d\theta}\right|_{\theta=\pi/4}=$ \rule{2cm}{0.4pt} (rounded off to one decimal place).
\item Let $C_n$ denote a cyclic group having $n$ elements. If there is a surjective group homomorphism from $C_n$ to $C_{30}$, then the total number of such distinct surjective homomorphisms is \rule{2cm}{0.4pt}.
\item $\displaystyle\int_0^3(|x-1|-x[x])\,dx=$ \rule{2cm}{0.4pt} (rounded off to one decimal place).
\item Let $A\in M_3(\mathbb C)$. Suppose the column vector $v=\begin{pmatrix}\sqrt5\,i\\2i\\x\end{pmatrix}$ in $\mathbb C^3$ belongs to the intersection of $\operatorname{nullspace}(A)$ and $\operatorname{rangespace}(A^T)$. Then $|x|=$ \rule{2cm}{0.4pt} (rounded off to one decimal place).
\item Let $P$ be a $5\times5$ real matrix with $\det(P)=2$. Let $Q$ be the matrix of cofactors of $P$. Then $\det(Q)=$ \rule{2cm}{0.4pt} (rounded off to one decimal place).
\item Let $\alpha,\beta$ and $\gamma$ be fixed real numbers such that
\[y(x)=C_1e^{-x}+C_2e^{2x}+\alpha xe^{-x},\quad\text{where }C_1\text{ and }C_2\text{ are arbitrary real constants},\]
is the general solution of the differential equation
\[\frac{d^2y}{dx^2}+\beta\frac{dy}{dx}+\gamma y=-e^{-x}.\]
Then $\alpha(\beta+\gamma)=$ \rule{2cm}{0.4pt} (rounded off to one decimal place).
\item Let the function $f:\mathbb R^2\to\mathbb R$ be defined by
\[f(x,y)=\begin{cases}
4x^2\tan^{-1}\left(\dfrac{y}{2x}\right)+y^3\tan^{-1}\left(\dfrac{x}{4y^2}\right),&xy\ne0\\
0,&\text{otherwise}.
\end{cases}\]
Then $\left(\dfrac{\partial}{\partial y}\left(\dfrac{\partial f}{\partial x}\right)\right)$ at $(0,0)$ is \rule{2cm}{0.4pt} (rounded off to two decimal places).
\item The double integral of $f(x,y)=x$ over the triangular region with vertices at $\left(-\frac12,\frac12\right),(1,2)$ and $(1,-1)$ is \rule{2cm}{0.4pt} (rounded off to one decimal place).
\item Let
\[f(x)=\int_0^{e^{\sin x}}\int_0^{\log_e y} e^{-(1/\sqrt2)}(1-t^2)\,dt\,dy.\]
Then $f'(\pi/4)=$ \rule{2cm}{0.4pt} (rounded off to one decimal place).
\item The volume of the tetrahedron bounded by the planes $x=1,y=2,z=3$ and $12x+8y+6z=70$ is \rule{2cm}{0.4pt} (rounded off to one decimal place).
\item Let $M\in M_3(\mathbb R)$. If
\[M\begin{pmatrix}\sqrt{15}\\\sqrt{15}\\0\end{pmatrix}=\begin{pmatrix}\sqrt3\\\sqrt2\\-5\end{pmatrix},\quad
M\begin{pmatrix}0\\0\\\sqrt{30}\end{pmatrix}=\begin{pmatrix}\sqrt{15}\\\sqrt{10}\\\sqrt5\end{pmatrix},\quad
M\begin{pmatrix}-\sqrt{15}\\\sqrt{15}\\0\end{pmatrix}=\begin{pmatrix}\sqrt{12}\\-\sqrt{18}\\0\end{pmatrix},\]
then $\det(M)=$ \rule{2cm}{0.4pt} (rounded off to one decimal place).
\item Let $a,b\in\mathbb R$. If the system of linear equations
\[\begin{aligned}
x+2y+2z&=1\\
2x+3y+z&=2\\
ax+5y+bz&=b
\end{aligned}\]
has infinitely many solutions, then $a+b=$ \rule{2cm}{0.4pt} (rounded off to one decimal place).
\item A fruit shop has $4$ different types of bananas. The number of ways in which $12$ bananas can be bought with at least one banana from each type, is \rule{2cm}{0.4pt}.
\item Let $A$ be a $3\times3$ real matrix such that given any column vector $x\in\mathbb R^3$, the column vector $Ax$ is the reflection of $x$ about the plane
\[\{(a,b,-a-b):a,b\in\mathbb R\}.\]
Then the sum of the diagonal elements of $A$ is \rule{2cm}{0.4pt} (rounded off to one decimal place).
\item Let $\binom nr$ denote the number of ways of choosing $r$ distinct objects out of $n$ distinct objects. Then,
\[\binom50+\binom61+\binom72+\binom83+\binom94+\binom{10}5+\binom{11}6=\rule{2cm}{0.4pt}.\]
\end{enumerate}
